Should a fixed inference budget buy more programs or more executions of
one program? Automated agent design makes the first option increasingly
accessible: a builder writes a population of LLM harnesses, each
specifying prompts, tools, and control flow around a frozen solver
\citep{adas2024,tthe2026,selfharness2026}. The appeal is that different
programs might handle different tasks better. To justify that
interpretation, however, a population's gains must be separated from
the extra attempts used to evaluate it.

The difficulty is that different observed successes do not require
different abilities. Two independent executions of a program with
success probability $0.5$ cover a task with probability $0.75$.
No program specialization has been introduced. Pass@$k$ measures
coverage over multiple samples; self-consistency and sampling-and-voting
ensembles aggregate multiple attempts into one answer
\citep{chen2021codex,selfconsistency2022,moreagents2024}.
An oracle that selects a correct answer after seeing all outcomes can
therefore show a population gain even when every member has the same
code. Averaging a few repeats does not automatically remove this
ambiguity: choosing the largest estimated success rate also favors
sampling noise.

\begin{figure}[t]
\centering
\setlength{\fboxsep}{8pt}
\fbox{\begin{minipage}[t][100pt][t]{0.445\linewidth}
\small
\textbf{One program, more executions}\par\medskip
\centering
$\boxed{H_0^{(1)}}\qquad\boxed{H_0^{(2)}}$\par\smallskip
Same code; independent errors\par\medskip
$p=0.5\quad\longrightarrow\quad1-(1-p)^2=0.75$\par\medskip
\raggedright
\textbf{Coverage increases by $25$\,\%.}\par
Program specialization remains zero.
\end{minipage}}\hfill
\fbox{\begin{minipage}[t][100pt][t]{0.445\linewidth}
\small
\textbf{A measured same-code reference}\par\medskip
\centering
MATH-500: nine identical \bare{} slots\par\smallskip
Three independent repeats per slot\par\medskip
$H_{\rm stable}=0\qquad\widehat H=\WpVthreeClonePlugin\,\%$\par\medskip
\raggedright
\textbf{The observed estimate is positive.}\par
Descriptive plugin on $383$ complete tasks.
\end{minipage}}
\caption{\textbf{More coverage does not identify more specialization.}
Left: a probability calculation for two independent executions of one
program. Right: the empirical clone arm produces a positive
repeat-averaged oracle estimate despite zero program specialization
(complete-task estimate; \S\ref{sec:f3-clones}).
The panels illustrate different quantities, not matched effect sizes.
Our control design compares eight generated programs plus \bare{} with
nine identical \bare{} slots, then separates repeat validation from
feature-based pre-execution selection.}
\label{fig:sampling-control}
\end{figure}


We ask what evidence distinguishes this coverage opportunity from
task specialization. Three questions organize the study. Did the
recorded executions contain additional correct answers? Does a member's
advantage on a task persist in a fresh execution? Can a rule choose
usefully before seeing the answer, and improve on spending the same
resources on a fixed program? These questions concern different
quantities. A positive answer to the first need not answer the other two.

We first report single-run coverage in the full population of $35$
generated MATH-500 programs, then test the specialization interpretation
on a selected subset. The central control holds population size fixed
and changes whether the programs differ: eight generated harnesses plus
\bare{} are compared with nine byte-identical baseline slots. Three
fresh repeats let us discover task-wise choices on two executions and
score them on the third. The same-code arm measures what this procedure
reports without program specialization. On the same panel, a selector frozen on
development tasks then tests selection from question features, and an
execution-matched oracle replay compares available coverage. Finally,
BIRD traces examine whether code differences actually reach execution,
providing a complementary diagnosis within the SQL setting.

The evidence yields three contributions. \textbf{First}, identical
programs produce positive repeat-averaged oracle headroom, so that
signal alone cannot identify program specialization. \textbf{Second},
generated programs have substantially more repeatable scored-outcome
patterns than clones, dominated by persistent losses: losses span
$\RcvLossTasks$ tasks ($\RcvLossPairs$ member--task pairs), whereas
persistent wins occur on one task. Stable differences therefore need
not provide useful task advantages; the clone-adjusted held-out gain is
$\RcvD$ pp. \textbf{Third}, the frozen selector gains
$\RcvPolicyGain$ pp, and the 27-execution oracle replay ties at
$\RcvReplayFullCoverage\%$. On this panel, more programs produce
neither a gain for the tested selector nor greater full-budget oracle
coverage than more executions of the baseline.
